% ### Erste Seite der Fluidkinematik ###
\centering
\vspace{188mm}
\section*{\myul{Fluidkinematik}}

\begin{tikzpicture}[remember picture, overlay]

% --- ERSTE REIHE: Am Seitenursprung anfangen
\setlength{\rowheight}{16mm}

    % X1) Eigene Zeile für die Überschrift
    \Tile{X1}{([xshift=10mm, yshift=-223mm]current page.north west)}{190mm}{8mm}
        {
            \titlebf{Fluidbewegung}
        }{0}{0}{0}{0}

    % A1) Erste Kachel unten anbauen
    \setlength{\cellwidth}{105mm}
    \Tile{A1}{(X1.south west)}{\cellwidth}{\rowheight}
        {
            \vspace{-\TileSep}
            \titlerm{Langrange}\\ \vspace{-3mm}
            \begin{equation*}
                \frac{\text{D} f}{\text{D} t} = \frac{\partial f}{\partial t} + u\,\frac{\partial f}{\partial x} + v\,\frac{\partial f}{\partial y} + w\,\frac{\partial f}{\partial z} = \frac{\partial f}{\partial t} + \underline{v}\cdot\nabla(f)
            \end{equation*}
        }{0}{1}{0}{0}

    % B1) Erste Kachel unten anbauen
    \setlength{\cellwidth}{85mm}
    \Tile{B1}{(A1.north east)}{\cellwidth}{\rowheight}
        {
            \vspace{-\TileSep}
            \titlerm{Euler}\\ \vspace{-3mm}
            \begin{equation*}
                \frac{\text{d} f}{\text{d} t} = \frac{\partial f}{\partial t} + \frac{\partial f}{\partial x}\,\frac{\text{d} x}{\text{d} t} + \frac{\partial f}{\partial y}\,\frac{\text{d} y}{\text{d} t} + \frac{\partial f}{\partial z}\,\frac{\text{d} z}{\text{d} t}
            \end{equation*}
        }{0}{0}{0}{0}

\end{tikzpicture}